% !PRINTED 41
\weberpage{41}
Beide Glieder sind einander gleich; denn sie gehen in einander über
durch Vertauschung der gleichwerthigen Symbole $a$ und $b$, oder $c$
und $d$. Ihre Summe kann also durch das doppelte erste Glied ersetzt
werden und somit ist diese Invariante $(\Delta,\ \Delta)^2$ auch dargestellt durch
\[
(\Delta,\ \Delta)^2 = (ab)^2(cd)^2(ac)(bd) = (ab)^2(cd)^2(ad)(bc).
\]
\textit{Beispiel 2.} Wir stellen uns ferner die Aufgabe, eine durch ein
symbolisches Product dargestellte Covariante oder Invariante einer
Form $f$ in den Wurzeldifferenzen dieser Form zu berechnen. Die
Methode will ich an einem Beispiele entwickeln. Es sei gegeben die
Form $f = a_x^4$; sie besitzt eine Invariante $i = (ab)^4$. Um sie in den
Wurzeln $\alpha_1$, $\alpha_2$, $\alpha_3$, $\alpha_4$ von $f$ darzustellen, bilden wir zunächst diese
Invariante durch vierte Ueberschiebung von $a_x^4$ mit:
\[
f = \alpha_{1x}\ \alpha_{2x}\ \alpha_{3x}\ \alpha_{4x}
\]
und erhalten nach Nr. 38 (I):
\weq{1}{i = (ab)^4 = (a\alpha_1)(a\alpha_2)(a\alpha_3)(a\alpha_4);}
auf diese Weise ist bereits ein Symbol $b$ durch die Wurzeln $\alpha_1\alpha_2\alpha_3\alpha_4$
ersetzt. Es handelt sich also noch darum, auch für das Symbol $a$ die
Wurzeln $\alpha_i$ einzuführen. Wir ersetzen zu dem Zwecke in (1) $a$ durch
$z$ und überschieben die Polare $b_z^4 = \alpha_{1z}\alpha_{2z}\alpha_{3z}\alpha_{4z}$ viermal über die
Polare $a_z^4 = \alpha_{1z}\alpha_{2z}\alpha_{3z}\alpha_{4z}$. So erhalten wir nach Nr. 38 (I) und An\-merkung:
\[
(a_z^4,\ b_z^4)^4 = (ab)^4 = (\alpha_{1z}\cdot\alpha_{2z}\cdot\alpha_{3z}\cdot\alpha_{4z},\ \alpha_{1z}\alpha_{2z}\alpha_{3z}\alpha_{4z})^4
\]
\[
= \frac{1}{4!}\sum (\alpha_{i_1}\alpha_1)(\alpha_{i_2}\alpha_2)(\alpha_{i_3}\alpha_3)(\alpha_{i_4}\alpha_4),
\]
oder mit Vernachlässigung jener Glieder der Summe, welche wegen
Factoren $(\alpha_{i_1}\alpha_{i_1}) = 0$ identisch verschwinden:
\[
(ab)^4 = i = \frac{1}{4!}\left\{(\alpha_1\alpha_2)(\alpha_3\alpha_4) + (\alpha_1\alpha_3)(\alpha_2\alpha_4) + (\alpha_1\alpha_4)(\alpha_2\alpha_3)\right\}^2.
\]
Ganz in derselben Weise verfährt man auch allgemein. Man ersetzt
zunächst ein Symbol der Covariante durch $y$, und bildet die
$n^{\text{te}}$ Ueberschiebung dieser Polare mit der Polare $f_{y^n} = \alpha_{1y}\alpha_{2y}\cdots\alpha_{ny}$.
So erhält man ein Aggregat von Gliedern, welche ein Symbol
weniger enthalten. In jedem dieser Glieder ersetzt man wiederum ein
Symbol durch $z$, und überschiebt dasselbe abermals $n$ mal über
$f_{z^n} = \alpha_{1z}\alpha_{2z}\cdots\alpha_{nz}$ etc. Dadurch ist die Zahl der Symbole um 2
reducirt, und so setzt man den Process fort.] % [sic] stray closing bracket as printed; no opening bracket visible on this page

40. \textit{Darstellung der Ueberschiebung durch Ueberschiebung von
Polaren. Erster Fall.} Bei Berechnung einer Ueberschiebung, welche
an zwei Producten von Formen ausgeführt werden soll, kann man
sich nun
%CONT>>


% ===== pdf idx 56, printed 42.0 =====
% !PRINTED 42
\weberpage{42}
%<<CONT
aber auch wieder des bereits in Nr. 34 bewiesenen Satzes bedienen.
Die Sätze über Eigenschaften der Polaren lassen sich alsdann auf die
Ueberschiebung übertragen. Ich will mich indess bei den folgenden
Entwicklungen auf den Fall beschränken, in dem $F$ und $\Phi$ nur als
Producte von je zwei Formen vorausgesetzt werden, da die Verallge\-
meinerung nur eine öftere Anwendung der Operationen des speciellen
Falles erheischt, und keine neuen Ueberlegungen nöthig macht. Ueber\-
dies sei zunächst vorausgesetzt, dass der Grad der beiden Producte
in $x$ gleich dem Grade $k$ der Ueberschiebung sei, dass also, wenn
\[
F = f\cdot\varphi = a_x^m\cdot b_x^n
\]
\[
\Phi = \psi\cdot\chi = \alpha_x^\mu\cdot\beta_x^\nu
\]
für den Exponenten der Ueberschiebung
\[
U = (F,\ \Phi)^k
\]
die Beziehung bestehe
\[
k = m+n = \mu+\nu.
\]
Unter dieser Voraussetzung besteht nun aber die $k^{\text{te}}$ Polare von
$F$ aus dem einzigen Gliede
\[
F_{y^k} = a_y^m\cdot b_y^n
\]
und ebenso
\[
\Phi_{y^k} = \alpha_y^\mu\cdot\beta_y^\nu.
\]
Faltet man das Product $F_{y^k}\cdot \Phi_{y^k}$ auf alle möglichen Arten, etwa
indem man die linearen Factoren der einen Polare einen Moment durch
$r_{ix}$, die der andern durch $\varrho_{kx}$ ersetzt und nun die einzige Combination
der $r_{ix}$ mit allen Permutationen der $\varrho_{ix}$ durch Faltung verbindet, so
bildet die Summe aller dieser Glieder, dividirt durch ihre Anzahl, die
$k^{\text{te}}$ Ueberschiebung von $F$ über $\Phi$. Ersetzt man wieder $r_{ix}$ und $\varrho_{ix}$
in geeigneter Weise durch die ursprünglichen Factoren $a_x$ $b_x$ $\alpha_x$ $\beta_x$,
fasst hierauf alle gleichen Glieder zu einem einzigen Gliede $L_s$ zu\-
sammen, so ist die Ueberschiebung dargestellt durch
\[
(F,\ \Phi)^k = \sum q_s L_s,
\]
wo $L_s$ von der Form $(a\alpha)^{\nu_1}(b\beta)^{\nu_2}(a\beta)^{\nu_3}(b\beta)^{\nu_4}$, $\nu_1+\nu_2+\nu_3+\nu_4=k$,
und $q_s$ constante Coefficienten sind, deren Summe gleich Eins ist.

41. \textit{Zweiter Fall.} Ist nun der Grad des Productes $F$, und der
des Productes $\Phi$ grösser als der der Ueberschiebung, oder höchstens
eines der beiden Producte von gleichem Grade wie die Ueberschiebung,
so ist die $k^{\text{te}}$ Polare von $F$ (vergl. Nr. 22)
\[
F_{y^k} = \binom{m}{0}\binom{n}{k} a_x^m b_x^{n-k} b_y^k + \binom{m}{1}\binom{n}{k-1} a_x^{m-1} a_y b_x^{n-k+1} b_y^{k-1} + \cdots = \sum c_i G_i
\]
und ebenso
%CONT>>


% ===== pdf idx 57, printed 43.0 =====
% !PRINTED 43
\weberpage{43}
%<<CONT
\[
\Phi_{y^k} = \binom{\mu}{0}\binom{\nu}{k} \alpha_x^\mu \beta_x^{\nu-k} \beta_y^k + \binom{\mu}{1}\binom{\nu}{k-1} \alpha_x^{\mu-1} \alpha_y \beta_x^{\nu-k+1} \beta_y^{k-1} + \cdots = \sum p_\varrho H_\varrho.
\]
Die Ueberschiebung erhält man, indem man das Product
\[
F_{y^k}\cdot \Phi_{y^k} = \sum c_i G_i \cdot \sum p_\varrho H_\varrho
\]
auf alle möglichen Arten nach den in $y$ linearen Factoren faltet. Sie
ist also dargestellt durch die Summe
\weq{I}{(F,\ \Phi)^k = \sum\sum c_i\cdot p_\varrho (G_i,\ H_\varrho)^k,}
wo nun das Glied $(G_i,\ H_\varrho)^k$ entsteht, indem man auf irgend zwei
Glieder der $k^{\text{ten}}$ Polaren $k$ mal den Faltungsprocess anwendet. Die Ueber\-
schiebung der beiden Formen $F$ und $\Phi$ ist also zurückgeführt auf
eine Summe von Ueberschiebungen solcher Formen $G_i$ und $H_\varrho$, deren
Grad in $y$ genau mit dem Grade der Ueberschiebung übereinstimmt.
Jede solche Ueberschiebung ist nach Nr. 40 dargestellt durch
\[
(G_i,\ H_\varrho)^k = \sum q_s L_s.
\]
Trägt man also dieselben in (I) ein, so erhält man:
\[
(F,\ \Phi)^k = \sum\sum c_i p_\varrho \sum q_s L_s = \sum\sum\sum c_i p_\varrho q_s L_s
\]
\[
= \sum \lambda_{i\varrho s}(a\alpha)^{\nu_1}(b\beta)^{\nu_2}(a\beta)^{\nu_3}(b\beta)^{\nu_4}\, a_x^{m-(\nu_1+\nu_3)}\, b_x^{n-(\nu_2+\nu_4)}\, \alpha_x^{\mu-(\nu_1+\nu_4)}\, \beta_x^{\nu-(\nu_2+\nu_3)}.
\]
Hierbei ist die Summe der Coefficienten $\lambda_{i\varrho s}$ gleich Eins. Denn
es ist
\[
\sum c_i = 1, \qquad \sum p_\varrho = 1 \quad \text{nach Nr. 23}
\]
und ebenso
\[
\sum q_s = 1 \quad \text{nach Nr. 40,}
\]
folglich auch
\[
\sum c_i \cdot \sum p_\varrho \cdot \sum q_s = \sum \lambda_{i\varrho s} = 1.
\]

42. \textit{Benachbarte Glieder einer Ueberschiebung.} Wie bei den Po\-
laren so spricht man auch hier von benachbarten Gliedern. Doch ist
die Mannigfaltigkeit der Nachbarschaft nun eine grössere geworden.
So können zwei Glieder benachbart sein, wenn in dem einen an
Stelle von
\weq{1}{a_x (b\alpha)\ \text{sich befindet}\ b_x (a\alpha),}
oder an Stelle von
\weq{2}{\alpha_x (a\beta)\ \text{sich befindet}\ \beta_x (a\alpha),}
während alle übrigen Factoren übereinstimmen. Im ersten Falle hat
die Differenz der betreffenden Glieder den Factor
%CONT>>


% ===== pdf idx 58, printed 44.0 =====
% !PRINTED 44
\weberpage{44}
%<<CONT
\[
a_x(b\alpha) - b_x(a\alpha) = a_x(ab), % checked against the GDZ scan (27 Oct 2026): printed italic a_x, overwritten in ink by a handwritten alpha (a reader's correction); the stroke under the = sign is also handwritten
\]
im zweiten Falle den Factor
\[
\alpha_x(a\beta) - \beta_x(a\alpha) = a_x(\alpha\beta),
\]
wie aus dem Identitätssatz Nr. 11 hervorgeht.
Die Differenz zweier in dieser Weise benachbarter Glieder hat
demnach immer einen der beiden Factoren
\[
(ab)\quad\text{oder}\quad(\alpha\beta).
\]
Die Nachbarschaft zweier Glieder kann aber auch lediglich durch
Klammerfactoren bedingt sein. Es kann die Stelle von
\weq{3}{(a\alpha)(b\beta)\ \text{das Product}\ (a\beta)(b\alpha)}
einnehmen; dann hat nach dem Identitätssatze die Differenz dieser
beiden Glieder
\[
G_1 - G_2 = \bar{G}\left\{(a\alpha)(b\beta) - (a\beta)(b\alpha)\right\}
\]
\[
= \bar{G}(ab)(\alpha\beta)
\]
den Factor $(ab)(\alpha\beta)$.

\glqq Diese Differenzen haben also in allen drei Fällen mindestens
einen Klammerfactor, der durch Faltung zwischen Formen $F$
allein, oder $\Phi$ allein entsteht, und nicht durch gegenseitige
Faltung von $F$ mit $\Phi$.\grqq{}
[So ist z.\,B. in der schon mehrfach berechneten Ueberschiebung
\[
U = \left((ab)^2 a_x b_x,\ (cd)^2 c_x d_x\right)^2
\]
\[
= \frac{(ab)^2(cd)^2}{2}\left\{(ac)(bd) + (ad)(bc)\right\},
\]
die Differenz der beiden Glieder dargestellt durch
\[
G_1 - G_2 = (ab)^2(cd)^2\left\{(ac)(bd) - (ad)(bc)\right\}
\]
\[
= (ab)^3(cd)^3.
\]
Die rechte Seite verschwindet aus bekannten Gründen, wenn $a$,
$b$, $c$, $d$ Symbole der cubischen Form $f = a_x^3$ sind. Dann ist also
\[
G_1 = G_2,
\]
folglich
\[
U = \frac{1}{2}\cdot 2 G_1 = (ab)^2(cd)^2(ac)(bd),
\]
was wir schon auf anderem Wege gezeigt haben. (Vergl. Nr. 39.)
Die eben besprochene Eigenschaft der Differenz benachbarter
Glieder kann man auch für die Differenz zweier beliebiger Glieder
$L_i$ und $L_p$ nachweisen. Denn es ist
\[
L_i - L_p = (L_i - L_{i+1}) + (L_{i+1} - L_{i+2}) + \cdots + (L_{p-1} - L_p).
\]


% ===== pdf idx 59, printed 45.0 =====
% !PRINTED 45
\weberpage{45}
Jede Differenz rechts besitzt den Factor $(ab)$ oder $(\alpha\beta)$, also ist
die Differenz links durch ein Aggregat von Gliedern darstellbar, welche
entweder den Factor $(ab)$ oder den Factor $(\alpha\beta)$ oder beide besitzen.

43. \textit{Anwendungen.} Man kann von diesen Sätzen interessante Anwendungen
machen, die für uns in der Folge von Wichtigkeit sein werden. Denken wir uns
eine Covariante $A_i$ der Form $f_n=a_x^n$, so können wir uns genau dieselbe
Covariante $B_i$ der Form $\varphi=\alpha_x^m$ bilden, wenn $m\geqq n$, indem
wir an der Potenz $\varphi^\varrho$ die nämlichen Faltungen vornehmen, die wir
an $f^\varrho$ machten, um die Form $A_i$ zu erhalten. Nennen wir der
Einfachheit halber $A_i$ das Modell, $B_i$ das Bild. Ist insbesondere
\[
\varphi=(ab)^{2\mu}a_x^{n-\mu}b_x^{n-\mu}=(aa_1)^{2\mu}a_x^{n-\mu}a_{1x}^{n-\mu}, % checked against the GDZ scan (27 Oct 2026): all four exponents are printed n-\mu (misprint for n-2\mu); the thin 2s written over the minus signs are handwritten
\]
also eine gerade Ueberschiebung der Form $f$ über sich selbst, so können wir
mit Hilfe der Eigenschaft benachbarter Glieder für Bilder $B_i$, welche aus
dieser Form nach den Modellen $A_i$ gestaltet werden, eine Reihe von Sätzen
ableiten.

In diesem Falle nämlich kann man aus der Potenz
\[
\varphi^\varrho=(aa_1)^{2\mu}a_x^{n-\mu}a_{1x}^{n-\mu}\cdot(bb_1)^{2\mu}b_x^{n-\mu}b_{1x}^{n-\mu}\cdots(rr_1)^{2\mu}r_x^{n-\mu}r_{1x}^{n-\mu} % checked against the GDZ scan (27 Oct 2026): exponents printed n-\mu; the marks over them are handwritten 2s, as above
\]
eine Menge Bilder $B_i$ nach dem Modelle $A_i$ herstellen, je nachdem man die
durch $A_i$ dictirten Faltungen der Potenz
\[
f^\varrho=a_x^n\cdot b_x^n\cdot c_x^n\cdots r_x^n
\]
in $\varphi^\varrho$ ausführt, d.\,h. je nachdem man in $\varphi^\varrho$ die
entsprechenden Faltungen an $a_x^{n-\mu}b_x^{n-\mu}$, oder an
$a_x^{n-\mu}b_{1x}^{n-\mu}$, oder an $a_{1x}^{n-\mu}b_x^{n-\mu}$, oder
endlich an $a_{1x}^{n-\mu}b_{1x}^{n-\mu}$ vornimmt.

Wir nennen dann wieder zwei Bilder $B_i$ und $B_i'$ benachbarte Bilder, wenn
sie dadurch aus einander hervorgehen, dass man in zwei symbolischen Factoren
derselben $a$ mit $a_1$ vertauscht. Ihre Differenz hat dann immer um einen
Factor $(aa_1)=(bb_1)=(cc_1)$ mehr als das einzelne Bild.

Ist z.\,B.
\[
\begin{aligned}
B_i&=(ab)(a_1b_1)G_i\\
B_i'&=(a_1b)(ab_1)G_i,
\end{aligned}
\]
dann ist:
\[
\begin{aligned}
B_i-B_i'&=G_i\{(ab)(a_1b_1)-(a_1b)(ab_1)\}\\
&=G_i(aa_1)(bb_1);
\end{aligned}
\]
oder ist:
\[
\begin{aligned}
B_i&=a_x(a_1b)G_i\\
B_i'&=a_{1x}(ab)G_i,
\end{aligned}
\]
so ist
\[
B_i-B_i'=G_i\{a_x(a_1b)-a_{1x}(ab)\}=G_i(aa_1)b_x.
\]


% ===== pdf idx 60, printed 46.0 =====
% !PRINTED 46
\weberpage{46}
Genau wie vorhin erkennt man daher, dass auch die Differenz zwei beliebiger
Bilder $B_i$ und $\overline{B}_i$ sich durch ein Aggregat von Gliedern
darstellen lässt, deren jedes einen Factor $(aa_1)$ mehr besitzt als das
einzelne Bild. Da nun ein Bild mindestens $(aa_1)^{2\mu}$ bereits zum Factor
hat, so erkennt man die Richtigkeit der Behauptung:

\glqq Die Differenz zweier Bilder $B_i$ ist gleich einem Aggregat von
Gliedern, welche den Klammerfactor $(ab)^{2\mu+1}$ oder wegen Nr.\,12
$(ab)^{2\mu+2}$ besitzen.\grqq{}

Man kann auch sagen:

\glqq Die Differenz zweier Bilder ist null, bis auf Glieder mit dem Factor
$(ab)^{2\mu+2}$.\grqq{}

Wir werden einen derartigen aus Klammerfactoren zusammengesetzten
charakteristischen Factor Modul nennen, und demnach diesen Satz in die
algebraische Form kleiden:
\weq{1}{\overline{B}_i-B_i\equiv 0\ \mathrm{mod}\ (ab)^{2\mu+2}.}
Hat man nun drei symbolische Producte $A$, so dass
\[
\begin{aligned}
A_1&=a_x(bc)\psi\\
A_2&=b_x(ca)\psi\\
A_3&=c_x(ab)\psi,
\end{aligned}
\]
dann verschwindet wegen des Identitätssatzes ihre Summe identisch. Das
Gleiche gilt für die erste Gattung der oben erwähnten vier Bilder,
\[
\begin{aligned}
B_1&=a_x(bc)\psi'\\
B_2&=b_x(ca)\psi'\\
B_3&=c_x(ab)\psi',
\end{aligned}
\]
die ja neben andern Klammerfactoren genau dieselben enthalten wie die
Modelle $A_i$.

Nimmt man demnach drei beliebige Bilder $\overline{B}_1\ \overline{B}_2\
\overline{B}_3$, so ist wegen Gleichung (1)
\[
\sum_1^3\overline{B}_i-\sum_1^3 B_i\equiv 0\ \mathrm{mod}\ (ab)^{2\mu+2},
\]
weil aber $\sum B_i=0$, so erhält man
\weq{2}{\sum_1^3\overline{B}_i\equiv 0\ \mathrm{mod}\ (ab)^{2\mu+2}.}
Dieses Resultat lässt uns aber einen sehr wichtigen Satz beweisen:

\glqq Ist $F(A_i)$ eine ganze Function von symbolischen Producten $A_i$ der
Art, dass weder ihre Glieder einzeln null sind, noch
%CONT>>


% ===== pdf idx 61, printed 47.0 =====
% !PRINTED 47
\weberpage{47}
%<<CONT
sich gegenseitig aufheben; so ist, wenn diese Function $F$ gleichwohl
identisch verschwindet, die entsprechende Function von Bildern, also von
symbolischen Producten $B_i$ der Form $\varphi=(f,f)^{2\mu}$
\[
F(\overline{B}_i)\equiv 0\ \mathrm{mod}\ (ab)^{2\mu+2}.
\]
\grqq{} Denn man kann setzen (vergl.\ \S\,1 Nr.\,12, und \S\,10 Nr.\,117)
\[
F(A_i)=0=\sum\{(bc)a_x+(ca)b_x+(ab)c_x\}\varphi
\]
und demnach für die genau entsprechenden Bilder (erster Gattung)
\[
F(B_i)=\sum\{(bc)a_x+(ca)b_x+(ab)c_x\}\varphi'=0.
\]
Sind alsdann $\overline{B}_i$ beliebige Bilder, so ist
\[
F(\overline{B}_i)-F(B_i)\equiv 0\ \mathrm{mod}\ (ab)^{2\mu+2};
\]
also weil der Subtrahend links verschwindet
\weq{3}{F(\overline{B}_i)\equiv 0\ \mathrm{mod}\ (ab)^{2\mu+2}.}
Wir werden später von diesem Satze eine hervorragende Anwendung machen.

44. \textit{Differenz der Ueberschiebung selbst mit einem ihrer Glieder.}
Die Eigenschaft, die wir für die Differenz zweier beliebiger Glieder einer
Ueberschiebung nachgewiesen haben, besitzt auch die Differenz zwischen einem
Gliede und der ganzen Ueberschiebung. Denn da in der Ueberschiebung die
Summe aller Coefficienten Eins ist, so folgt durch Subtraction der beiden
Gleichungen
\[
U=\sum\lambda_{i\varrho s}L_s
\]
\[
L_q=L_q\sum\lambda_{i\varrho s}
\]
die Gleichung
\[
L_q=U+\sum_{i,\varrho,s}\lambda_{i\varrho s}(L_q-L_s)
\]
oder, wenn wir links in den Differenzen die Factoren $(ab)$ resp.
$(\alpha\beta)$ absondern:
\weq{I}{L_q=U+(ab)^{\varepsilon_1}(\alpha\beta)^{\varepsilon_2}\sum\lambda_{i\varrho s}P_s,\quad (\varepsilon_1\ \mathrm{u.}\ \varepsilon_2=0\ \mathrm{oder}\ 1)\ \varepsilon_1\geqq\varepsilon_2=0.}
Nun sind die symbolischen Producte $(ab)^{\varepsilon_1}(\alpha\beta)^{\varepsilon_2}P_s$
Glieder, welche irgend einen Klammerfactor, wie ihn die Ueberschiebung $U$
involviren würde, weniger besitzen, dafür aber mindestens einen mehr, welcher
dieser Ueberschiebung $U$ fremd ist, insofern er nur durch Faltung der Formen
von $F$ unter sich, oder der Formen $\Phi$ unter sich entstehen kann. Daraus
lernen wir zweierlei: Erstens, diese Glieder in den
%CONT>>


% ===== pdf idx 62, printed 48.0 =====
% !PRINTED 48
\weberpage{48}
%<<CONT
Summen rechts gehören Ueberschiebungen an, die einen um mindestens Eine
Einheit niedrigeren Grad $k-1$ haben als die Ueberschiebung $U$; zweitens,
diese Ueberschiebungen beziehen sich auf Formen niedrigeren Grades, welche,
sei es aus $F$ oder sei es aus $\Phi$, durch Faltung entstanden sind.

Wir gewinnen also zunächst den Satz:

\glqq Jedes Glied einer Ueberschiebung lässt sich darstellen durch die
Ueberschiebung selbst plus einer Summe von Gliedern, welche niedrigeren
Ueberschiebungen von Formen niedrigeren Grades angehören.\grqq{}

Ersetzen wir nun vermöge dieses Satzes abermals die Glieder
\[
(ab)^{\varepsilon_1}(\alpha\beta)^{\varepsilon_2}P_i
\]
durch ihre Ueberschiebungen und fahren in diesem Processe so lange fort, bis
wir zu nullten Ueberschiebungen von Formen gelangen, die aus $F$ oder $\Phi$
durch vielfache Faltung entstanden sind, so werden wir endlich zu dem Satze
geführt:

\glqq Jedes Glied einer Ueberschiebung lässt sich darstellen durch eine
Summe von Ueberschiebungen,\grqq{} oder
\weq{II}{L=U+\sum c_i^{(1)}U_i^{(1)}+\sum c_i^{(2)}U_i^{(2)}+\cdots+\sum c_i^{(k)}U_i^{(k)}.}
Dabei sind die Grössen $c_i^{(\varrho)}$ Constante, die Grössen
$U_i^{(\varrho)}$ der Reihe nach Ueberschiebungen $(k-1)^{\mathrm{ten}}$,
$(k-2)^{\mathrm{ten}}$ bis nullten Grades von Formen, die aus $F$ oder
$\Phi$ durch einmalige, zweimalige resp. $k$ malige Faltung entstanden sind.

Diese beiden Sätze (I) und (II) entsprechen den analogen Sätzen (I) und (II)
Nr.\,25 über Polaren.

45. \textit{Jedes symbolische Product lässt sich auf eine Summe einfacher
Ueberschiebungen zweier Formen reduciren.} Der letzte Satz kann noch
erweitert werden, indem wir ihn auf ein beliebiges symbolisches Product,
ohne Rücksicht darauf, ob ihm die Bedeutung einer wirklichen binären Form
zukommt oder nicht, ausdehnen. Das Princip, das uns die Darstellung eines
solchen symbolischen Productes durch eine Summe von einfachen
Ueberschiebungen ermöglicht, liegt in der wiederholten Einführung neuer
Symbole, die so lange fortgesetzt wird, bis in jedem Gliede der
entstehenden Summe nur mehr zwei Symbole übrig sind. (Vgl.\ auch Camille
Jordan, Liouville Journal, sér.\,3, Bd.\,2 und 5.) Ich will der
Deutlichkeit halber ein Beispiel voraussenden.

Es sei gegeben das symbolische Product:
\[
P=(ab)^2(ac)^2(bc)^2 a_x b_x c_x.
\]


% ===== pdf idx 63, printed 49.0 =====
% !PRINTED 49
\weberpage{49}
Halten wir in demselben die beiden Symbole $a$ und $b$ fest, ersetzen dagegen das Symbol $c$ durch eine Variable $y$, so wird aus $P$, abgesehen vom Factor $(xy)$, der Ausdruck:
\[
G_\varrho = (ab)^2\,a_y^2\,b_y^2\,a_x\,b_x.
\]
Derselbe stellt nichts anderes dar, als ein Glied der vierten Polare einer Form:
\[
\alpha_x^6 = (ab)^2\,a_x^3\,b_x^3.
\]
Dieses Glied lässt sich in die Polare $(\alpha_x^6)_{y^4}$ plus einer Summe von Polaren niedrigerer Formen entwickeln, die aus $\alpha_x^6$ durch Faltung entstehen. Das Resultat der Entwicklung können wir direct aus dem Beispiel Nr. 26 ableiten, indem wir in der Gleichung
\[
a_x^2\,b_x^2\,a_y\,b_y = F_{y^2} - \frac{2}{5}\,(xy)\,\Phi_y
\]
$x$ mit $y$ vertauschen und mit $(ab)^2$ multipliciren; dies liefert
\[
(ab)^2\,a_y^2\,b_y^2\,a_x\,b_x = (\alpha_x^6)_{y^4} + \frac{2}{5}\,(xy)\,(\beta_x^4)_{y^3},
\]
wobei mit $\beta_x^4$ die Form $(ab)^3\,a_x^2\,b_x^2$ bezeichnet ist. Ersetzen wir hierin $y_1$ wieder durch $c_2$ und $y_2$ durch $-c_1$ und multipliciren mit dem aus dem unterdrückten Factor $(xy)$ durch diese Substitution wieder hergestellten Factor $c_x$, so erhalten wir:
\[
P = (ab)^2\,(ac)^2\,(bc)^2\,a_x\,b_x\,c_x = (\alpha_x^6,\, c_x^5)^4 + \frac{2}{5}\,(\beta_x^4,\, c_x^5)^3.
\]
Das symbolische Product $P$ ist somit durch eine vierte und eine dritte Ueberschiebung dargestellt.

46. \textit{Allgemeine Darstellung.} Um also ein symbolisches Product (mit beliebig vielen Reihen von Symbolen)
\[
\begin{array}{c}
a,\ b,\ c \dots \text{etc.} \\
\alpha,\ \beta,\ \gamma \dots \text{etc.} \\
\mathfrak{a},\ \mathfrak{b},\ \mathfrak{c} \dots \text{etc.} \\
\cdot\ \ \cdot\ \ \cdot\ \ \cdot\ \ \cdot\ \ \cdot\ \ \cdot
\end{array}
\]
durch eine Summe von Ueberschiebungen darzustellen, verfahren wir folgendermassen. Wir halten eine Reihe von Symbolen fest, etwa die Symbole $a, b, c \dots$ etc., ersetzen alle übrigen durch Variable $y, z, u, v \dots$ etc. und sondern einen Augenblick die dadurch aus Factoren erster Art entstehenden identischen Covarianten $(xy), (xz), (xu), (xv) \dots$ ab. Der so entstehende Ausdruck ist dann nichts anderes, als ein Polarenglied einer Form, die wir mit $f$ bezeichnen wollen. Dieses Glied lässt sich aber nach Nr. 24 und 31 in eine Reihe von Polaren entwickeln, welche Formen angehören mögen, die wir mit
\[
\varphi = A_x^u, \quad \psi = B_x^v, \quad \chi = C_x^\varrho \dots \text{etc.}
\]
%CONT>>


% ===== pdf idx 64, printed 50.0 =====
% !PRINTED 50
\weberpage{50}
%<<CONT
bezeichnen. Ersetzt man nun in dieser Reihe die Formen $\varphi, \psi, \chi \dots$ etc. durch das ihnen zugeordnete Symbol $A, B, C$ etc., substituirt sodann für die Variabeln $y, z, u \dots$ wieder die ursprünglichen Symbole
\[
\begin{array}{c}
\alpha,\ \beta,\ \gamma,\ \dots \text{etc.} \\
\mathfrak{a},\ \mathfrak{b},\ \mathfrak{c},\ \dots \text{etc.}
\end{array}
\]
und multiplicirt nun mit den vorhin abgesonderten Factoren $(xy), (xz) \dots$ d. i. mit:
\[
\begin{array}{c}
\alpha_x,\ \beta_x,\ \gamma_x \dots \\
\mathfrak{a}_x,\ \mathfrak{b}_x,\ \mathfrak{c}_x \dots,
\end{array}
\]
so entsteht auf der linken Seite das ursprüngliche Product $P$. Rechts dagegen erhalten wir eine Summe von symbolischen Producten, deren jedes statt der ganzen Symbolreihe
\[
a,\ b,\ c,\ \dots
\]
nur je ein einziges enthält, sei es das Symbol $A$ der Form $\varphi$, oder das Symbol $B$ der Form $\psi, \dots$ etc. Ein solches Glied enthält also weniger Symbole wie vorher, und wir können nun in jedem derselben, genau auf dieselbe Weise wie in $P$ selbst, die Zahl der Symbole reduciren.

Auf diese Weise gelangen wir schliesslich zu einer Summe von Gliedern, deren jedes nur mehr zwei Symbole enthält. Solche symbolische Producte haben wir aber als einfache Ueberschiebungen definirt, und wir sind somit berechtigt den Satz aufzustellen: \glqq\gesperrt{Jedes symbolische Product lässt sich durch eine Summe von einfachen Ueberschiebungen darstellen.}\grqq{}

\S~4. \textbf{Die einfachsten Ueberschiebungen.}

47. \textit{Die Functionaldeterminante.} Unter den Ueberschiebungen haben insbesondere zwei eine hervorragende Bedeutung. Es sind gerade die beiden einfachsten Ueberschiebungen zweier Formen, nämlich die erste und die zweite. Die erste ist bekannt unter dem Namen \glqq Functionaldeterminante\grqq{} oder auch \glqq Jacobi'sche Determinante\grqq{} (vergl. Bd. I \S~9). Sie ist, wenn
\[
f = a_x^m, \quad \varphi = b_x^n
\]
die beiden Formen sind, dargestellt durch
\[
(f, \varphi) = a_x^{m-1}\,b_x^{n-1}\,(ab),
\]
und unterscheidet sich von der im ersten Bande \S~9 definirten Function nur um den Factor $m \cdot n$. Man erkennt dies unmittelbar aus der unsymbolischen Schreibweise dieser Form (vgl. Nr. 35, III)
%CONT>>


% ===== pdf idx 65, printed 51.0 =====
% !PRINTED 51
\weberpage{51}
%<<CONT
\weq{I}{(f, \varphi) = \begin{vmatrix} f_1, & f_2 \\ \varphi_1, & \varphi_2 \end{vmatrix},}
wobei
\[
f_i = \frac{1}{m}\,\frac{\partial f}{\partial x_i}, \quad \varphi_i = \frac{1}{n}\cdot\frac{\partial \varphi}{\partial x_i}.
\]
Diese Determinante lässt sich übrigens, was für die folgenden Betrachtungen bequemer ist, auch als eine dreigliedrige darstellen. Das symbolische Product
\[
(ab)\,a_x\,b_x = (a_1b_2 - b_1a_2)(a_1x_1 + a_2x_2)(b_1x_1 + b_2x_2)
\]
hat nämlich, wie die Determinantentheorie lehrt (vergl. Bd. I Nr. 40, Anmerkung), als Determinante geschrieben die Matrix
\[
\begin{vmatrix} a_1^2, & a_1a_2, & a_2^2 \\ b_1^2, & b_1b_2, & b_2^2 \\ x_2^2, & -x_1x_2, & x_2^2 \end{vmatrix} = (ab)\,a_x\,b_x.
\]
Multiplicirt man hier die erste Zeile mit $a_x^{m-2}$, die zweite mit $b_x^{n-2}$ und berücksichtigt, dass
\[
\begin{array}{ll}
a_1^2\,a_x^{m-2} = f_{11}, & b_1^2\,b_x^{n-2} = \varphi_{11} \\
a_1a_2\,a_x^{m-2} = f_{12}, & b_1b_2\,b_x^{n-2} = \varphi_{12} \\
a_2^2\,a_x^{m-2} = f_{22}, & b_2^2\,b_x^{n-2} = \varphi_{22},
\end{array}
\]
so erhält man für die Functionaldeterminante die dreigliedrige Determinante:
\weq{II}{(f, \varphi) = (ab)\,a_x^{m-1}\,b_x^{n-1} = \begin{vmatrix} f_{11}, & f_{12}, & f_{22} \\ \varphi_{11}, & \varphi_{12}, & \varphi_{22} \\ x_2^2, & -x_1x_2, & x_1^2 \end{vmatrix}.}

48. \textit{Die Functionaldeterminante als Form ungeraden Charakters.} Da die Functionaldeterminante nur einen einzigen Klammerfactor besitzt, so ist sie eine Form ungeraden Charakters (\textit{forme gauche}). Wir können zeigen, dass ihr Quadrat sich als eine ganze Function von Formen geraden Charakters darstellen lassen muss, eine Eigenschaft, die sie mit allen schiefen Co- und Invarianten theilt. (Vergl. Nr. 8.)

In der That, vertauschen wir in (II) Nr. 47 die erste und letzte Colonne und multipliciren die zweite Colonne mit $-2$, so erhalten wir
\[
+2\,(f, \varphi) = \begin{vmatrix} f_{22}, & -2f_{12}, & f_{11} \\ \varphi_{22}, & -2\varphi_{12}, & \varphi_{11} \\ x_2^2, & +2x_1x_2, & x_1^2 \end{vmatrix}.
\]


% ===== pdf idx 66, printed 52.0 =====
% !PRINTED 52
\weberpage{52}
Bilden wir nun das Product dieser Determinante mit der ursprünglichen, so kommt:
\[
\begin{aligned}
2\,(f, \varphi)\cdot(f, \varphi) &= \begin{vmatrix} f_{11}, & f_{12}, & f_{22} \\ \varphi_{11}, & \varphi_{12}, & \varphi_{22} \\ x_1^2, & -x_1x_2, & x_2^2 \end{vmatrix} \cdot \begin{vmatrix} f_{22}, & -2f_{12}, & f_{11} \\ \varphi_{22}, & -2\varphi_{12}, & \varphi_{11} \\ x_2^2, & 2x_1x_2, & x_1^2 \end{vmatrix} \\
&= \begin{vmatrix} 2(f_{11}f_{22} - f_{12}^2), & f_{22}\varphi_{11} - 2f_{12}\varphi_{12} + f_{11}\varphi_{22}, & f \\ f_{22}\varphi_{11} - 2f_{12}\varphi_{12} + f_{11}\varphi_{22}, & 2(\varphi_{22}\varphi_{11} - \varphi_{12}^2), & \varphi \\ f, & \varphi, & 0 \end{vmatrix}.
\end{aligned}
\]
Hierin ist nach Nr. 45, III:
\[
\begin{array}{l}
f_{11}\varphi_{22} - 2f_{12}\varphi_{12} + f_{22}\varphi_{11} = (f, \varphi)^2 \\
2(f_{11}f_{22} - f_{12}^2) = (f, f)^2 \\
2(\varphi_{11}\varphi_{22} - \varphi_{12}^2) = (\varphi, \varphi)^2.
\end{array}
\]
Demnach erhalten wir als Ausdruck für das Quadrat der Functionaldeterminante $J = (f, \varphi)$
\[
J^2 = \frac{1}{2} \begin{vmatrix} (f, f)^2, & (f, \varphi)^2, & f \\ (f, \varphi)^2, & (\varphi, \varphi)^2, & \varphi \\ f, & \varphi, & 0 \end{vmatrix}
\]
oder
\weq{I}{J^2 = -\frac{1}{2}\left\{(f, f)^2\cdot\varphi^2 - 2(f, \varphi)^2 f\cdot\varphi + (\varphi, \varphi)^2\cdot f^2\right\}.}

Die Formen $f$ und $\varphi$ sowohl, als ihre zweiten Ueberschiebungen, sind aber Formen geraden Charakters; denn $f$ und $\varphi$ haben keine und ihre zweiten Ueberschiebungen haben zwei Klammerfactoren, gehen also bei Transformation bis auf eine gerade Potenz des Moduls $\Delta$ in sich selbst über.

Genau in derselben Weise kann man auch das Product zweier verschiedener Functionaldeterminanten darstellen. Man erhält
\[
\begin{aligned}
2\,(f, \varphi)\cdot(\chi, \psi) &= \begin{vmatrix} f_{11}, & f_{12}, & f_{22} \\ \varphi_{11}, & \varphi_{12}, & \varphi_{22} \\ x_1^2, & -x_1x_2, & x_2^2 \end{vmatrix} \cdot \begin{vmatrix} \chi_{22}, & -2\chi_{12}, & \chi_{11} \\ \psi_{22}, & -2\psi_{12}, & \psi_{11} \\ x_2^2, & +2x_1x_2, & x_1^2 \end{vmatrix} \\
&= \begin{vmatrix} (f, \chi)^2, & (\varphi, \chi)^2, & \chi \\ (f, \psi)^2, & (\varphi, \psi)^2, & \psi \\ f, & \varphi, & 0 \end{vmatrix}
\end{aligned}
\]
oder
\weq{II}{(f, \varphi)\cdot(\chi, \psi) = -\frac{1}{2}\left\{(f, \chi)^2\varphi\cdot\psi - (f, \psi)^2\varphi\cdot\chi - (\varphi, \chi)^2 f\cdot\psi + (\varphi, \psi)^2 f\cdot\chi\right\}.}


% ===== pdf idx 67, printed 53.0 =====
% !PRINTED 53
\weberpage{53}

49. \textit{Die Functionaldeterminante der zweiten Polaren dreier Formen.} Bildet man von den drei Formen
\[
f = a_x^m, \quad \varphi = b_x^n, \quad \psi = c_x^p
\]
die zweiten Polaren
\[
\begin{aligned}
f_{y^2} &= f_{11}y_1{}^2 + 2f_{12}y_1y_2 + f_{22}y_2{}^2 \\
\varphi_{y^2} &= \varphi_{11}y_1{}^2 + 2\varphi_{12}y_1y_2 + \varphi_{22}y_2{}^2 \\
\psi_{y^2} &= \psi_{11}y_1{}^2 + 2\psi_{12}y_1y_2 + \psi_{22}y_2{}^2,
\end{aligned}
\]
so können wir im Sylvester'schen Sinne (vergl. Bd. I. Nr. 143) die Resultante der drei in $y$ quadratischen Formen als Functionaldeterminante betrachten und erhalten, abgesehen von einem Factor 2, als Ausdruck derselben die Determinante
\weq{I}{
\begin{aligned}
U &= \begin{vmatrix} f_{11}, & f_{12}, & f_{22} \\ \varphi_{11}, & \varphi_{12}, & \varphi_{22} \\ \psi_{11}, & \psi_{12}, & \psi_{22} \end{vmatrix}
= a_x^{m-2}b_x^{n-2}c_x^{p-2} \begin{vmatrix} a_1{}^2, & a_1a_2, & a_2{}^2 \\ b_1{}^2, & b_1b_2, & b_2{}^2 \\ c_1{}^2, & c_1c_2, & c_2{}^2 \end{vmatrix} \\
&= (ab)(ac)(bc)\,a_x^{m-2}b_x^{n-2}c_x^{p-2}.
\end{aligned}
}
Dieselbe Determinante geht auch aus $J = (f, \varphi)$ Nr. 47 (II) hervor, wenn wir $x_1$ durch $-\psi_2$ und $x_2$ durch $\psi_1$ ersetzen.

Auch diese Covariante $U$ ist eine Form ungeraden Charakters und wir können sie ganz in derselben Weise wie die Functionaldeterminante $J$ durch gerade Formen darstellen. Multipliciren wir sie nämlich mit
\[
2\,U = \begin{vmatrix} f_{22}, & -2f_{12}, & f_{11} \\ \varphi_{22}, & -2\varphi_{12}, & \varphi_{11} \\ \psi_{22}, & -2\psi_{12}, & \psi_{11} \end{vmatrix},
\]
so erhalten wir:
% checked against the GDZ scan (27 Oct 2026): a round dot stands in the left margin beside this display; not part of the text
\weq{II}{
2\,U^2 = \begin{vmatrix} (f,\,f)^2, & (f,\,\varphi)^2, & (f,\,\psi)^2 \\ (f,\,\varphi)^2, & (\varphi,\,\varphi)^2, & (\varphi,\,\psi)^2 \\ (f,\,\psi)^2, & (\varphi,\,\psi)^2, & (\psi,\,\psi)^2 \end{vmatrix}.
}
Die Determinante rechts ist eine Covariante, die, wie wir später sehen werden, unter die Gattung der Combinanten zu rechnen ist. Ihre Elemente sind die Coefficienten von $\lambda$ in der zweiten Ueberschiebung
\[
(F,\,F)^2
\]
der Form
\[
F = \lambda_1 f + \lambda_2 \varphi + \lambda_3 \psi,
\]
über sich selbst. Denn bilden wir die zweite Polare
\[
F_{y^2} = \lambda_1 f_{y^2} + \lambda_2 \varphi_{y^2} + \lambda_3 \psi_{y^2},
\]
%CONT>>


% ===== pdf idx 68, printed 54.0 =====
% !PRINTED 54
\weberpage{54}

%<<CONT
so erhält man daraus die Ueberschiebung, wenn man in dem Producte
\[
\{\lambda_1 f_{y^2} + \lambda_2 \varphi_{y^2} + \lambda_3 \psi_{y^2}\}\{\lambda_1 f_{y^2} + \lambda_2 \varphi_{y^2} + \lambda_3 \psi_{y^2}\}
\]
alle nach der Multiplication entstehenden Glieder nach $y$ zweimal faltet. Dies ist, da wir hier einfache Formen zweiten Grades in $y$ haben, immer nur auf eine einzige Art möglich, so dass wir erhalten:
\[
(F,\,F)^2 = \lambda_1{}^2(f,\,f)^2 + \lambda_2{}^2(\varphi,\,\varphi)^2 + \cdots + 2\lambda_1\lambda_2(f,\,\varphi)^2 + \cdots
\]
50. \textit{Relation zwischen drei und vier binären Formen.} Das in diesem Paragraph nun schon zum wiederholten Male benutzte Princip, durch Determinantenmultiplication Relationen zu erhalten, wollen wir noch benutzen, um Relationen zwischen beliebigen Formen herzustellen.

Wir erhalten eine Relation zwischen den drei beliebigen Formen $f, \varphi$ und $\psi$, wenn wir die beiden Determinanten vom Werthe null
\weq{I}{
D = \begin{vmatrix} f_{11}, & f_{12}, & f_{22}, & 0 \\ \varphi_{11}, & \varphi_{12}, & \varphi_{22}, & 0 \\ \psi_{11}, & \psi_{12}, & \psi_{22}, & 0 \\ x_1{}^2, & -x_1x_2, & x_2{}^2, & 0 \end{vmatrix}
}
und
\weq{II}{
2\,D = \begin{vmatrix} f_{22}, & -2f_{12}, & f_{11}, & 0 \\ \varphi_{22}, & -2\varphi_{12}, & \varphi_{11}, & 0 \\ \psi_{22}, & -2\psi_{12}, & \psi_{11}, & 0 \\ x_2{}^2, & +2x_1x_2, & x_1{}^2, & 0 \end{vmatrix}
}
mit einander multipliciren. Das Resultat der Multiplication ist:
% checked against the GDZ scan (27 Oct 2026): no mark in the left margin; a stray dot stands to the right of this display, below the number (III); not part of the text
\weq{III}{
0 = \begin{vmatrix} (f,\,f)^2, & (f,\,\varphi)^2, & (f,\,\psi)^2, & f \\ (f,\,\varphi)^2, & (\varphi,\,\varphi)^2, & (\varphi,\,\psi)^2, & \varphi \\ (f,\,\psi)^2, & (\varphi,\,\psi)^2, & (\psi,\,\psi)^2, & \psi \\ f, & \varphi, & \psi, & 0 \end{vmatrix} = \varDelta.
}
Nach den Elementen der letzten Zeile und Reihe ausgewerthet, liefert sie, wenn man die Unterdeterminanten zweiter Ordnung mit $A_{\substack{4k\\4i}}$ bezeichnet (vergl. Bd. I. Nr. 64 und Beispiel 65)
\[
\begin{aligned}
\varDelta &= -\sum_{4i} a_{4i}\,a_{k4}\,A_{\substack{4k\\4i}} \\
&= -\left\{A_{\substack{41\\41}}f^2 + A_{\substack{42\\42}}\varphi^2 + A_{\substack{43\\43}}\psi^2 + 2A_{\substack{41\\42}}f\cdot\varphi + 2A_{\substack{41\\43}}f\cdot\psi + 2A_{\substack{42\\43}}\varphi\cdot\psi\right\},
\end{aligned}
\]
oder in leicht verständlicher Abkürzung:
\[
\varDelta = -(A_1f + A_2\varphi + A_3\psi)^2 = 0.
\]


% ===== pdf idx 69, printed 55.0 =====
% !PRINTED 55
\weberpage{55}

Vertauscht man in den beiden Determinanten (I) und (II) wieder $x_1$ mit $\chi_2$ und $x_2$ mit $-\chi_1$, so erhält man durch Multiplication derselben die Relation:
\[
\begin{vmatrix}
(f,\,f)^2, & (f,\,\varphi)^2, & (f,\,\psi)^2, & (f,\,\chi)^2 \\
(f,\,\varphi)^2, & (\varphi,\,\varphi)^2, & (\varphi,\,\psi)^2, & (\varphi,\,\chi)^2 \\
(f,\,\psi)^2, & (\psi,\,\varphi)^2, & (\psi,\,\psi)^2, & (\psi,\,\chi)^2 \\
(f,\,\chi)^2, & (\chi,\,\varphi)^2, & (\chi,\,\psi)^2, & (\chi,\,\chi)^2
\end{vmatrix} = 0,
\]
und diese Determinante giebt eine Relation zwischen Covarianten von vier beliebigen Formen. Wir werden später, insbesondere bei der Theorie der quadratischen Formen, von den bisher entwickelten Relationen vielfach Gebrauch machen.

51. \textit{Die Functionaldeterminante $(f, \varphi)$ und die erste Polare $(f\cdot\varphi)_y$.} Die Functionaldeterminante giebt auch noch in anderer Richtung zu Relationen Veranlassung, die besonders für die symbolische Rechnung von Nutzen sind. Hierher gehört ihre Beziehung zur Polare $(f\cdot\varphi)_y$. Diese Polare ist nämlich, wenn $f = a_x^m = b_x^m\cdots$ und $\varphi = \alpha_x^n = \beta_x^n\cdots$, dargestellt durch:
\weq{1}{
(f\cdot\varphi)_y = \frac{m}{m+n}\cdot a_x^{m-1}a_y\alpha_x^n + \frac{n}{m+n}a_x^m\alpha_x^{n-1}\alpha_y = c_0G_0 + c_1G_1 \text{ (vgl. Nr. 22).}
}
Die Differenz der beiden Glieder ist:
\[
G_1 - G_2 = a_x^{m-1}\alpha_x^{n-1}(a_y\alpha_x - \alpha_ya_x) = a_x^{m-1}\alpha_x^{n-1}(a\alpha)(yx),
\]
oder
\weq{2}{
(yx)\cdot(f,\varphi) = a_x^{m-1}a_y\alpha_x^n - a_x^m\alpha_x^{n-1}\alpha_y.
}
Eliminiren wir aus beiden Gleichungen (1) und (2) einmal das erste, sodann das zweite Glied rechts, so erhalten wir:
\weq{I}{
\begin{aligned}
(f\cdot\varphi)_y - \frac{m}{m+n}(yx)\cdot(f,\varphi) &= a_x^m\alpha_x^{n-1}\alpha_y = f\cdot\varphi_y \\
(f\cdot\varphi)_y + \frac{n}{m+n}(yx)\cdot(f,\varphi) &= a_x^{m-1}a_y\alpha_x^n = f_y\cdot\varphi
\end{aligned}
}
Diese beiden Gleichungen, von denen die zweite auch durch Vertauschung von $f$ mit $\varphi$ aus der ersten erhalten werden kann, setzen die Polare $(f\cdot\varphi)_y$ mit der Functionaldeterminante $(f,\varphi)$ und den Polaren $f_y$ und $\varphi_y$ in Beziehung.

Dabei haben wir nicht vorausgesetzt, dass $f$ und $\varphi$, unsymbolisch genommen, wirkliche binäre Formen sind. Die Relationen (I) gelten also auch, wenn wir in ihnen $f$ durch $a_x^{m-1}$ und $\varphi$ durch $\alpha_x^{n-1}$ ersetzen und nachträglich mit $(a\alpha)$ multipliciren. Dann geht aber
\[
(f\cdot\varphi)_y \text{ über in } (a\alpha)\cdot(a_x^{m-1}\cdot\alpha_x^{n-1})_y = (f,\varphi)_y,
\]
%CONT>>


% ===== pdf idx 70, printed 56.0 =====
% !PRINTED 56
\weberpage{56}

%<<CONT
ferner:
\[
(f,\varphi) \text{ über in } (a\alpha)(a_x^{m-1},\,\alpha_x^{n-1}) = (a\alpha)^2a_x^{m-2}\alpha_x^{n-2} = (f,\varphi)^2.
\]
Demnach können wir aus den Relationen (I) die beiden neuen gewinnen:
\weq{II}{
\begin{aligned}
(f,\varphi)_y - \frac{m-1}{m+n-2}(yx)\cdot(f,\varphi)^2 &= (a\alpha)\,a_x^{m-1}\alpha_y\alpha_x^{n-2} \\
(f,\varphi)_y + \frac{n-1}{m+n-2}(yx)\cdot(f,\varphi)^2 &= (a\alpha)\,a_x^{m-2}a_y\alpha_x^{n-1}
\end{aligned}
}
In ihnen sind die Ausdrücke rechts die Glieder der Polare $(f,\varphi)_y$ und die Relationen (I) und (II) sind somit nichts anderes, als die in Nr. 25 gegebene Darstellung eines Gliedes durch seine Polare und die Polaren von niedrigeren Formen.

52. \textit{Die Functionaldeterminante von $(f,\varphi)$ und $f$ oder $\varphi$.} Eine besondere Rolle spielen bei späteren Untersuchungen die Functionaldeterminanten von Functionaldeterminanten, also die erste Ueberschiebung von $(f,\varphi)$ über eine dritte Form $\psi$, welche wir symbolisch bezeichnen mit
\[
((f,\varphi),\psi) = U.
\]
Diese Formen haben nämlich die Eigenschaft, dass sie sich immer durch ein Aggregat niedrigerer Formen darstellen lassen. Nehmen wir zunächst an, $\psi$ sei eine der beiden Formen $f$ oder $\varphi$ selbst. Die erste Ueberschiebung $U = ((f,\varphi),f)$ entsteht nämlich, wenn wir in der Polare $(f,\varphi)_y$ die Variabeln $y_1$ durch $b_2$, und $y_2$ durch $-b_1$ ersetzen (vergl. Nr. 35) und mit $b_x^{m-1}$ multipliciren. Aus der zweiten der Relationen (II) erhalten wir demnach:
\[
((f,\varphi),f) = (a\alpha)(ab)a_x^{m-2}b_x^{m-1}\alpha_x^{n-1} - \frac{n-1}{m+n-2}(f,\varphi)^2\cdot b_x\cdot b_x^{m-1}.
\]
Hierin ist das zweite Glied rechts bereits das Product von zwei einfacheren Formen $(f,\varphi)^2$ und $b_x^m = f$. Das erste Glied $G_1$ lässt sich mit Hilfe des Identitätssatzes in ein ähnliches Product verwandeln. Vertauscht man nämlich in ihm die gleichwerthigen Symbole $a$ und $b$, so wird aus dem Gliede
\[
G_1 = (a\alpha)(ab)a_x^{m-2}b_x^{m-1}\alpha_x^{n-1}
\]
wiederum
\[
G_1 = -(b\alpha)(ab)b_x^{m-2}a_x^{m-1}\alpha_x^{n-1}.
\]
Durch Addition erhält man:
\[
2G_1 = (ab)a_x^{m-2}b_x^{m-2}\alpha_x^{n-1}\{(a\alpha)b_x - (b\alpha)a_x\},
\]
oder, weil gemäss dem Identitätssatze die Differenz in der Klammer gleich $(ab)\alpha_x$ ist,
\[
G_1 = \frac{1}{2}(ab)^2a_x^{m-2}b_x^{m-2}\alpha_x^n = \frac{1}{2}(f,f)^2\cdot\varphi.
\]


% ===== pdf idx 71, printed 57.0 =====
% !PRINTED 57
\weberpage{57}
Die Functionaldeterminante $((f,\ \varphi),\ f)$ ist somit dargestellt durch:

\weq{I}{((f,\ \varphi),\ f) = \frac{1}{2}(f,f)^2\cdot\varphi - \frac{n-1}{m+n-2}\cdot(f,\ \varphi)^2\cdot f.}

Dabei ist natürlich vorausgesetzt, dass beide Formen $f$ und $\varphi$ von höherem als dem ersten Grade in $x$ sind.

Vertauschen wir in der Relation (I) $f$ mit $\varphi$, so erhalten wir auch die andere erste Ueberschiebung $((\varphi,\ f),\ \varphi) = -((f,\ \varphi),\ \varphi)$ dargestellt durch die Gleichung:

\weq{II}{((f,\ \varphi),\ \varphi) = -\frac{1}{2}(\varphi,\ \varphi)^2\cdot f + \frac{m-1}{m+n-2}(f,\ \varphi)^2\cdot\varphi\,.}

53. \textit{Die Functionaldeterminante von $(f,\ \varphi)$ und $\psi$.} Ist nun $\psi = r_x^p$ eine von $f$ und $\varphi$ verschiedene Form (sie kann auch selbst wieder eine Functionaldeterminante sein), so können wir die erste Ueberschiebung

\[
U = ((f,\ \varphi),\ \psi)
\]

auf folgende Weise durch Formen niederern Grades ausdrücken. Wir ersetzen zunächst wieder in der zweiten der beiden Relationen (II) Nr. 51 $y$ durch $r$, d. h. $y_1$ durch $r_2$, $y_2$ durch $-r_1$ und erhalten nach Multiplication mit $r_x^{n-1}$

\weq{1}{((f,\ \varphi),\ \psi) = (a\alpha)(ar)a_x^{m-2}\alpha_x^{n-1}r_x^{p-1} - \frac{n-1}{m+n-2}(f,\ \varphi)^2\cdot\psi.}

Zur Umformung des ersten Gliedes benutzen wir den Productsatz

\[
(a\alpha)(ar)\alpha_x r_x = \frac{1}{2}\left\{(a\alpha)^2 r_x^2 + (ar)_x^2\alpha_x^2 - (\alpha r)^2 a_x^2\right\}\cdot
\]

Substituiren wir diesen Werth von $(\alpha a)(ar)\alpha_x r_x$ in (1), so kommt:

\[
((f,\ \varphi),\ \psi) = \frac{1}{2}\left\{(a\alpha)^2 a_x^{m-2}\alpha_x^{n-2}\cdot\psi + (ar)^2 a_x^{m-2}r_x^{p-2}\cdot\varphi - (ar)^2\alpha_x^{m-2}r_x^{p-2}\cdot f\right\}
\]
\[
- \frac{n-1}{m+n-2}(f,\ \varphi)^2\cdot\psi
\]
\[
= \frac{1}{2}(f,\ \varphi)^2\cdot\psi + \frac{1}{2}(f,\ \psi)^2\cdot\varphi - \frac{1}{2}(\varphi,\ \psi)^2\cdot f - \frac{n-1}{m+n-2}(f,\ \varphi)^2\cdot\psi,
\]

oder endlich:

\weq{III}{((f,\ \varphi),\ \psi) = \frac{m-n}{2(m+n-2)}(f,\ \varphi)^2\cdot\psi + \frac{1}{2}\left\{(f,\ \psi)^2\cdot\varphi - (\varphi,\ \psi)^2\cdot f\right\}.}

Wir sind somit berechtigt, den Satz aufzustellen:

\glqq Die Functionaldeterminante einer Form $\psi$ mit einer andern Functionaldeterminante ist, wenn alle Formen vom höheren als vom ersten Grade in $x$ sind, immer durch Producte niedrigerer Formen ausdrückbar.\grqq{}

Aus der Relation (III) gehen selbstverständlich die beiden früheren Relationen (I) oder (II) hervor, wenn wir $\psi$ durch $f$ oder $\varphi$ ersetzen.
%CONT>>


% ===== pdf idx 72, printed 58.0 =====
% !PRINTED 58
\weberpage{58}
54. \textit{Die zweite Ueberschiebung einer Form $f$ über sich selbst.} Als zweites Beispiel einer einfachen Ueberschiebung wollen wir noch die Covariante $(ab)^2 a_x^{m-2}b_x^{m-2}$ betrachten. Diese Form wurde zuerst von Hesse als eine Covariante einer Form $f = a_x^m = b_x^m$ entdeckt (Crelle Journal Bd. 28) und führt nach ihm den Namen Hesse'sche Determinante oder Hesse'sche Form. Sie ist dargestellt symbolisch durch

\[
(f,\ f)^2 = (ab)^2 a_x^{m-2}b_x^{m-2} = H,
\]

unsymbolisch durch:

\[
H = \frac{2}{m^2(m-1)^2}
\begin{vmatrix}
\dfrac{\partial^2 f}{\partial x_1{}^2}, & \dfrac{\partial^2 f}{\partial x_1\,\partial x_2} \\[2ex]
\dfrac{\partial^2 f}{\partial x_1\,\partial x_2}, & \dfrac{\partial^2 f}{\partial x_2{}^2}
\end{vmatrix}\,.
\]

Man kann sie auch auffassen als Functionaldeterminante der beiden ersten Differentialquotienten $\dfrac{\partial f}{\partial x_1}$ und $\dfrac{\partial f}{\partial x_2}$. Wir haben diese Covariante bereits in mehreren Beispielen kennen gelernt. So ist die Discriminante der quadratischen Form $f = a_x^2 = b_x^2$ nichts anderes als die Hesse'sche Form

\[
H = (ab)^2 = 2(\bar a_0\bar a_2 - \bar a_1{}^2),
\]

ebenso war die in Nr. 37 berechnete Covariante zweiten Grades einer cubischen Form die Hesse'sche Covariante

\[
H = (ab)^2 a_x b_x = 2\left\{(\bar a_0\bar a_2 - \bar a_1{}^2)x_1{}^2 + (\bar a_0\bar a_3 - a_1\bar a_2)x_1 x_2 + \cdots\right\}.
\]

55. \textit{Lehrsatz von Hesse.} Von dieser Covariante hatte Hesse den allgemeinen Satz aufgestellt und bewiesen:

\glqq Ist die Determinante

\[
H =
\begin{vmatrix}
\dfrac{\partial^2 f}{\partial x_1{}^2}, & \dfrac{\partial^2 f}{\partial x_1\,\partial x_2}, & \dfrac{\partial^2 f}{\partial x_1\,\partial x_3} & \cdots & \dfrac{\partial^2 f}{\partial x_1\,\partial x_n} \\
\dfrac{\partial^2 f}{\partial x_1\,\partial x_2}, & \dfrac{\partial^2 f}{\partial x_2{}^2}, & \dfrac{\partial^2 f}{\partial x_2\,\partial x_3} & \cdots & \cdot \\
\cdot & \cdot & \cdot & \cdots & \cdot \\
\dfrac{\partial^2 f}{\partial x_1\,\partial x_n}, & \dfrac{\partial^2 f}{\partial x_2\,\partial x_n}, & \cdots & \cdots & \dfrac{\partial^2 f}{\partial x_n{}^2}
\end{vmatrix}
\]

der zweiten Differentialquotienten einer Function $f$ mit $n$ Variabeln identisch Null, so reducirt sich die Dimension der Form $f$ um einen Grad.\grqq{}

Im Jahre 1876 zeigten nun Gordan und Nöther, dass der Satz in dieser Allgemeinheit unrichtig ist.

Durch das Verschwinden der Hesse'schen Form reduciren sich zwar binäre Formen auf reine Potenzen eines linearen Ausdruckes, Formen mit drei Variabeln, also ternäre auf binäre, quaternäre auf ternäre; dagegen giebt es bereits quinäre Formen, deren $H$ zwar
%CONT>>


% ===== pdf idx 73, printed 59.0 =====
% !PRINTED 59
\weberpage{59}
%<<CONT
identisch verschwindet, die sich aber keineswegs auf quaternäre reduciren lassen.

Wir geben hier den Beweis des von Hesse aufgestellten Satzes für binäre Formen, indem wir ihn folgendermassen formuliren.

\glqq Ist die Covariante $(ab)^2 a_x^{m-2}b_x^{m-2}$ der Form $f(x_1 x_2)$ identisch null, so ist $f(x_1 x_2)$ eine reine Potenz: $(\bar a_1 x_1 + \bar a_2 x_2)^n$.\grqq{}

Zum Beweise bringen wir zunächst die Voraussetzung

\weq{1}{H = (ab)^2 a_x^{m-2}b_x^{m-2} = 0}

durch eine andere Relation zum Ausdruck. Wir multipliciren Gleichung (1) mit $\dfrac{(xy)^2}{2}$, und können dann schreiben:

\[
\frac{(xy)^2}{2}(ab)^2 a_x^{m-2}b_x^{m-2} = \frac{1}{2}(a_x b_y - b_x a_y)^2 a_x^{m-2}b_x^{m-2}
\]
\[
= \frac{1}{2}(a_x{}^2 b_x{}^2 - 2a_x b_y b_x a_y + b_x{}^2 a_y{}^2)a_x^{m-2}b_x^{m-2},
\]

oder, weil $a_x^m\cdot b_y^2 b_x^{n-2} = b_x^m a_y^2 a_x^{m-2} = f\cdot f_{y^2}$,

\weq{2}{\frac{(xy)^2}{2}H = f\cdot f_{y^2} - f_y{}^2 = 0.}

Zu dieser Voraussetzung (2) tritt noch eine zweite, dass nämlich jede Form $f(x)$ mindestens einen linearen Factor $\alpha_x$ besitzt (vgl. Fundamentalsatz der Algebra Bd. I \S~12), oder, um ganz allgemein zu sein, dass $f$ von der Form ist:

\weq{3}{f = \alpha_x^{\varrho}\cdot\psi_x^{\sigma} = \alpha_x^{\varrho}\cdot\psi\,,^{*)}\quad(\sigma + \varrho = m).}

Hier ist $\alpha_x$ wirklicher vielfacher Factor von $f$, $\varrho$ sein Exponent und $\psi$ das Product aller anderen Factoren. Unter diesen beiden Voraussetzungen lautet dann die Behauptung:

\[
\varrho = m\,.
\]

Wir bilden nun die erste und zweite Polare von $f = \alpha_x^{\varrho}\cdot\psi$, und erhalten nach Nr. 22:

\weq{4}{m\cdot f_y = \varrho\cdot\alpha_x^{\varrho-1}\alpha_y\cdot\psi + \sigma\cdot\alpha_x^{\varrho}\cdot\psi_x^{\sigma-1}\psi_y}

\weq{5}{m(m-1)f_{y^2} = \varrho(\varrho-1)\alpha_x^{\varrho-2}\alpha_y{}^2\cdot\psi + 2\varrho\cdot\sigma\alpha_x^{\varrho-1}\alpha_y\psi_x^{\sigma-1}\psi_y + \sigma(\sigma-1)\alpha_x^{\varrho}\cdot\psi_x^{\sigma-2}\psi_y\,.}

Quadriren wir also den Ausdruck (4) der ersten Polare und subtra-
%CONT>>

\medskip
\noindent{}$^{*)}$ Man darf annehmen $\varrho > 1$. Für $\varrho = 1$ wird der Beweis genau ebenso geführt.


% ===== pdf idx 74, printed 60.0 =====
% !PRINTED 60
\weberpage{60}
%<<CONT
hiren das $(m-1)$fache dieses Quadrates von dem $m$fachen Producte der Form $f$ in die Form (5), so erhalten wir:

\[
m^2(m-1)\{f\cdot f_{y^2} - f_y{}^2\} = \{m\varrho(\varrho-1) - (m-1)\varrho^2\}\alpha_x^{2\varrho-2}\alpha_y{}^2\cdot\psi^2
\]
\[
+ \{2m\varrho\sigma - 2\varrho\sigma(m-1)\}\alpha_x^{2\varrho-1}\alpha_y\cdot\psi_y\cdot\psi
\]
\[
+ \{m\sigma(\sigma-1)\psi\cdot\psi_{y^2} - (m-1)\sigma^2\psi_y{}^2\}\alpha_x^{2\varrho} = 0.
\]

Da nun in dieser Identität die späteren Glieder rechts den Factor $\alpha_x$ in einer höheren Potenz enthalten als das erste, so muss dessen Zahlencoëfficient für sich verschwinden, d. h. es muss sein:

\[
m\varrho(\varrho-1) = (m-1)\varrho^2
\]

oder

\[
m = \varrho\,,
\]

was zu beweisen war.

\begin{center}
\S~5. \textbf{Der Aronhold'sche Process.}
\end{center}

56. \textit{Aufgabe.} Viele Fragen der Algebra und Geometrie erheischen die Lösung folgender Aufgabe: Gegeben ist eine Form $f$ und eine ihrer Invarianten (Covarianten) $i$; gegeben ist ferner eine Form $\varphi$ von gleichem Grade in $x$ wie die Form $f$. Dann besitzt das lineare System $f + \lambda\varphi$ eine entsprechende simultane Invariante $J$, welche aus $i$ entsteht, wenn man darin die Coefficienten $\bar a_i$ von $f$ durch die Coefficienten $\bar a_i + \lambda\bar\alpha_i$ von $f + \lambda\varphi$ ersetzt$^{*)}$. Diese Invariante $J$ lässt sich nach Potenzen von $\lambda$ anordnen, und die Coefficienten der Potenzen von $\lambda$ werden alsdann selbst wieder simultane Invarianten von $f$ und $\varphi$ sein. Die Aufgabe nun, die wir uns stellen, ist: womöglich eine independente Darstellung der Coefficienten von $\lambda^k$ zu geben. Es wird sich zeigen, dass sich, so lange $f$ und $\varphi$ von einander unabhängig sind, immer ein Process angeben lässt, eng verwandt mit dem Polarenprocess, vermöge welchem jeder dieser Coefficienten unmittelbar aus der Invariante $i$ hergestellt werden kann; im anderen Falle erhält man für dieselben Recursionsformeln.

Um die Begriffe zu fixiren, nehmen wir an, es sei:

\[
f = a_x^n = b_x^n = \cdots \text{ etc.}
\]
\[
= \bar a_0 x_1^n + \binom{n}{1}\bar a_1 x_1^{n-1}x_2 + \binom{n}{2}\bar a_2 x_1^{n-2}x_2{}^2 + \cdots,
\]
\[
\varphi = \alpha_x^n = \beta_x^n = \cdots \text{ etc.}
\]
\[
= \bar\alpha_0 x^n + \binom{n}{1}\bar\alpha_1 x_1^{n-1}x_2 + \binom{n}{2}\bar\alpha_2 x_1^{n-2}x_2 + \cdots, % [sic]
\]

\medskip
\noindent{}$^{*)}$ Ist $\varphi$ eine Form höheren als $n^{\text{ten}}$ Grades, so bildet man die geeignete Polare und kann sich dann die gleiche Frage für die Form $f + \lambda\varphi_{yk}$ vorlegen.


% ===== pdf idx 75, printed 61.0 =====
